\documentclass[a4paper]{article}

\usepackage{graphicx}
\usepackage{amsmath}
\usepackage{amssymb}
\usepackage{float}
\usepackage{listings}
\usepackage{xcolor}
\usepackage{booktabs}
\usepackage[top=2cm, bottom=2cm, right=2.5cm, left=2.5cm]{geometry}

% پکیج‌های رسم نمودار (در صورت نیاز به رسم داخلی)
\usepackage{pgfplots}
\pgfplotsset{compat=1.17}

% تنظیمات نمایش کد برنامه
% تنظیمات نمایش کد برنامه
\lstset{
	frame=tb,
	language=Matlab,
	aboveskip=3mm,
	belowskip=3mm,
	showstringspaces=false,
	columns=flexible,
	basicstyle={\small\ttfamily},
	numbers=left,
	numberstyle=\tiny\color{gray},
	keywordstyle=\color{blue},
	commentstyle=\color{green!50!black},
	stringstyle=\color{red},
	breaklines=true,
	breakatwhitespace=true,
	tabsize=4
	% خط direction=LTR حذف شد
}
% بسته xepersian باید آخرین بسته‌ای باشد که فراخوانی می‌شود
\usepackage{xepersian}

% تنظیم فونت‌ها (اگر فونت وزیر را ندارید، نام فونت دیگری مثل XB Niloofar را جایگزین کنید)
\settextfont{Vazirmatn}
\setlatintextfont{Times New Roman}

\title{پروژه درس کنترل خطی\\ \large}
\author{مبین بهروان-حامد باطانی}
\date{\today}

\begin{document}
	
	\maketitle

	\section{انواع سیستم‌های تعلیق خودرو (بخش الف)}
	سیستم‌های تعلیق بر اساس توانایی آن‌ها در تغییر مشخصات دینامیکی و میزان مصرف انرژی به سه دسته کلی تقسیم می‌شوند:
	
	\subsection{سیستم تعلیق غیرفعال \lr{(Passive Suspension)}}
	\textbf{اصول عملکرد:} این سیستم‌ها از اجزای با مشخصات ثابت تشکیل شده‌اند. انرژی ناشی از ناهمواری‌های جاده توسط فنر ذخیره شده و توسط دمپر (کمک‌فنر) مستهلک می‌شود.
	\begin{itemize}
		\item \textbf{اجزای اصلی:} فنر \lr{(Spring)} و کمک‌فنر هیدرولیکی \lr{(Damper)} با ضریب میرایی ثابت.
		\item \textbf{کاربردها:} اکثر خودروهای سواری اقتصادی و معمولی.
		\item \textbf{مزایا:} سادگی طراحی، هزینه پایین، قابلیت اطمینان بالا و عدم نیاز به منبع انرژی خارجی.
		\item \textbf{محدودیت‌ها:} عملکرد آن محدود به فرکانس تشدید طبیعی سیستم است و نمی‌تواند همزمان نرمی سواری و پایداری را در شرایط متغیر جاده تأمین کند \lr{(Trade-off)}.
	\end{itemize}
	
	\subsection{سیستم تعلیق نیمه‌فعال \lr{(Semi-Active Suspension)}}
	\textbf{اصول عملکرد:} در این سیستم‌ها، ضریب میرایی کمک‌فنر می‌تواند به صورت بلادرنگ تغییر کند، اما سیستم نمی‌تواند به بدنه انرژی تزریق کند (تنها انرژی را تلف می‌کند).
	\begin{itemize}
		\item \textbf{اجزای اصلی:} دمپرهای مگنتو-رئولوژیکال \lr{(MR)} یا شیرهای الکترونیکی متغیر، سنسورهای شتاب و واحد کنترل الکترونیکی \lr{(ECU)}.
		\item \textbf{کاربردها:} خودروهای اسپرت و لوکس مدرن.
		\item \textbf{مزایا:} بهبود قابل توجه عملکرد نسبت به سیستم غیرفعال با مصرف انرژی اندک (در حد باتری خودرو).
		\item \textbf{محدودیت‌ها:} هزینه بالاتر و پیچیدگی کنترلی بیشتر نسبت به مدل غیرفعال.
	\end{itemize}
	
	\subsection{سیستم تعلیق تمام‌فعال \lr{(Fully Active Suspension)}}
	\textbf{اصول عملکرد:} این سیستم از یک عملگر \lr{(Actuator)} موازی با فنر استفاده می‌کند که می‌تواند به سیستم نیرو وارد کرده و بدنه خودرو را بلند یا پایین بیاورد.
	\begin{itemize}
		\item \textbf{اجزای اصلی:} پمپ‌های هیدرولیک فشار بالا یا موتورهای الکتریکی خطی، سنسورهای پیشرفته و پردازنده‌های سریع.
		\item \textbf{کاربردها:} خودروهای بسیار لوکس، نظامی و تحقیقاتی.
		\item \textbf{مزایا:} کنترل کامل بر تکان‌های بدنه، حذف کامل ارتعاشات و حفظ ارتفاع خودرو در پیچ‌ها.
		\item \textbf{محدودیت‌ها:} قیمت بسیار بالا، وزن زیاد، پیچیدگی فنی و مصرف انرژی قابل توجه.
	\end{itemize}
	
	\section{مدل‌های استاندارد دینامیکی (بخش ب)}
	برای تحلیل و طراحی کنترلر، مدل‌سازی ریاضی خودرو ضروری است. سه مدل استاندارد در ادبیات موضوع وجود دارد:
	
	\begin{table}[H]
		\centering
		\caption{مقایسه مدل‌های استاندارد تعلیق}
		\begin{tabular}{|c|c|p{8cm}|}
			\hline
			\textbf{نام مدل} & \textbf{درجات آزادی \lr{(DOF)}} & \textbf{توضیحات و کاربرد} \\
			\hline
			یک‌چهارم خودرو \lr{(Quarter-car)} & 2 & 
			شامل جرم بدنه ($m_s$) و جرم چرخ ($m_u$). مناسب برای تحلیل ارتعاشات عمودی و طراحی اولیه کنترلر. ساده‌ترین و پرکاربردترین مدل در آموزش است. \\
			\hline
			نیم‌خودرو \lr{(Half-car)} & 4 & 
			شامل حرکت عمودی مرکز جرم و حرکت پیچ \lr{(Pitch)} یا غلتش \lr{(Roll)}. برای بررسی اثرات ترمزگیری و شتاب‌گیری مناسب است. \\
			\hline
			خودروی کامل \lr{(Full-car)} & 7 & 
			شامل حرکت عمودی، پیچ و غلتش بدنه به همراه حرکت عمودی هر چهار چرخ. دقیق‌ترین مدل برای شبیه‌سازی واقعی و مانورهای پیچیده است. \\
			\hline
		\end{tabular}
	\end{table}
	
	\subsection{تحلیل مدل‌ها}
	\begin{itemize}
		\item \textbf{مدل \lr{Quarter-car}:} با وجود سادگی و عدم توانایی در مدل‌سازی اثرات هندسی خودرو، برای بررسی مفاهیم پایه ایزولاسیون ارتعاش بسیار کارآمد است.
		\item \textbf{مدل \lr{Full-car}:} اگرچه دقیق‌ترین است، اما به دلیل تعداد زیاد متغیرهای حالت و پیچیدگی معادلات، هزینه محاسباتی بالایی دارد و معمولاً در مراحل نهایی طراحی استفاده می‌شود.
	\end{itemize}
	
	\section{اهداف کنترلی و مصالحه طراحی (بخش ج)}
	در طراحی سیستم تعلیق، دو هدف اصلی وجود دارد که ذاتاً با یکدیگر در تضاد هستند:
	
	\subsection{1. آسایش سفر \lr{(Ride Comfort)}}
	این شاخص به میزان ارتعاشات منتقل شده به سرنشینان اشاره دارد. معمولاً با اندازه‌گیری \textbf{شتاب عمودی بدنه خودرو} ($\ddot{z}_s$) سنجیده می‌شود. هرچه دامنه و فرکانس این شتاب کمتر باشد، آسایش سفر بیشتر است.
	
	\subsection{2. خوش‌دستی جاده \lr{(Road Holding)}}
	این شاخص به توانایی تایرها در حفظ تماس دائمی با سطح جاده اشاره دارد. معیار سنجش آن، تغییرات \textbf{نیروی دینامیکی تایر} است. اگر نیروی تماس تایر با جاده کم شود، قابلیت فرمان‌پذیری و ترمزگیری خودرو به شدت کاهش می‌یابد.
	
	\subsection{تضاد \lr{(Trade-off)} بنیادی}
	بهبود یکی از این دو هدف، معمولاً منجر به تضعیف دیگری می‌شود:
	\begin{itemize}
		\item سیستم تعلیق \textbf{نرم} \lr{(Soft)}: آسایش عالی ایجاد می‌کند اما باعث نوسان زیاد بدنه و کاهش پایداری خودرو در پیچ‌ها می‌شود.
		\item سیستم تعلیق \textbf{سخت} \lr{(Hard)}: پایداری عالی دارد (مناسب خودروهای مسابقه‌ای) اما تمام ناهمواری‌ها را به سرنشین منتقل می‌کند.
\end{itemize}
	\subsection{انتخاب هدف طراحی پروژه}
	برای ادامه این پروژه، هدف \textbf{«آسایش سفر» \lr{(Ride Comfort)}} به عنوان هدف اصلی انتخاب می‌شود.
	\\
	\textbf{توجیه:} در خودروهای سواری شهری، راحتی سرنشین اولویت اول مشتریان است. علاوه بر این، کاهش شتاب بدنه به طور ملموس‌تری عملکرد کنترلر طراحی شده را در نمودارها نشان می‌دهد و برای یک پروژه آموزشی معیار مناسب‌تری است. البته همزمان قیودی برای جلوگیری از جدا شدن چرخ از سطح جاده \lr{(Road Holding)} در طراحی لحاظ خواهد شد.
	
	% ---------------------------------------------------------------------
	% شروع بخش دوم: مدل‌سازی ریاضی (Quarter-Car Model)
	% ---------------------------------------------------------------------
	
	\section{مدل‌سازی ریاضی و دینامیک سیستم (بخش‌های \lr{a} و \lr{b})}
	در این بخش، مدل یک‌چهارم خودرو \lr{(Quarter-Car Model)} که در شکل 1 صورت سوال ارائه شده است، تحلیل می‌شود. این مدل شامل دو جرم، فنرها و دمپر است که دینامیک ارتعاشات عمودی خودرو را شبیه‌سازی می‌کند.
	
	\subsection{الف) تحلیل نیروها و استخراج معادلات حرکت}
	برای بدست آوردن معادلات دیفرانسیل حاکم بر سیستم، از قانون دوم نیوتن ($\sum F = ma$) استفاده می‌کنیم. متغیرهای جابجایی $x_1$ و $x_2$ نسبت به وضعیت تعادل استاتیکی اندازه‌گیری می‌شوند، بنابراین نیروی وزن ($mg$) با نیروی اولیه فنرها خنثی شده و در معادلات ظاهر نمی‌شود.
	
	\subsubsection*{1. تحلیل جرم فنربندی شده \lr{(Sprung Mass, $M$)}}
	این جرم نماینده بدنه خودرو (شامل شاسی و سرنشینان) است. نیروهای وارد بر این جرم عبارتند از:
	\begin{itemize}
		\item \textbf{نیروی فنر تعلیق ($F_k$):} متناسب با اختلاف جابجایی نسبی دو جرم است. اگر بدنه ($x_1$) بیشتر از چرخ بالا برود، فنر کشیده شده و نیرو به سمت پایین وارد می‌کند: $-K_1(x_1 - x_2)$.
		\item \textbf{نیروی دمپر ($F_d$):} متناسب با سرعت نسبی دو جرم است و انرژی را مستهلک می‌کند: $-D(\dot{x}_1 - \dot{x}_2)$.
		\item \textbf{نیروی عملگر فعال ($u$):} نیروی کنترلی که توسط سیستم تعلیق فعال بین بدنه و چرخ اعمال می‌شود. جهت مثبت آن به سمت بالا فرض شده است.
	\end{itemize}
	
	بنابراین معادله حرکت نیوتن برای جرم $M$ به صورت زیر است:
	\begin{equation} \label{eq:M_dynamic}
		M \ddot{x}_1 = -K_1(x_1 - x_2) - D(\dot{x}_1 - \dot{x}_2) + u
	\end{equation}
	
	\subsubsection*{2. تحلیل جرم فنربندی نشده \lr{(Unsprung Mass, $m$)}}
	این جرم نماینده مجموعه چرخ، تایر و اکسل است. نیروهای وارد بر آن عبارتند از:
	\begin{itemize}
		\item \textbf{نیروهای عکس‌العمل تعلیق:} طبق قانون سوم نیوتن، نیروهای فنر و دمپر که به بدنه وارد می‌شوند، با جهت مخالف به چرخ وارد می‌گردند: $+K_1(x_1 - x_2) + D(\dot{x}_1 - \dot{x}_2)$.
		\item \textbf{نیروی عکس‌العمل عملگر:} نیروی $u$ با جهت منفی (به سمت پایین) به چرخ وارد می‌شود.
		\item \textbf{نیروی تایر ($F_{tire}$):} تایر مانند یک فنر با سختی $K_2$ مدل می‌شود. اگر جاده ($y_r$) بالاتر از موقعیت چرخ ($x_2$) باشد، تایر فشرده شده و نیرو به سمت بالا اعمال می‌کند: $+K_2(y_r - x_2)$.
	\end{itemize}
	
	معادله حرکت نیوتن برای جرم $m$ به صورت زیر خواهد بود:
	\begin{equation} \label{eq:m_dynamic}
		m \ddot{x}_2 = K_1(x_1 - x_2) + D(\dot{x}_1 - \dot{x}_2) - u + K_2(y_r - x_2)
	\end{equation}
	
	\subsection{ب) نمایش فضای حالت \lr{(State-Space Representation)}}
	برای طراحی کنترلر، معادلات دیفرانسیل مرتبه دوم بالا باید به فرم استاندارد فضای حالت مرتبه اول تبدیل شوند.
	
	\subsubsection*{1. تعریف متغیرهای حالت}
	بردار حالت $x$ شامل جابجایی‌ها و سرعت‌های جرم‌ها تعریف می‌شود:
	\begin{equation}
		x = \begin{bmatrix} x_1 \\ x_2 \\ v_1 \\ v_2 \end{bmatrix} = \begin{bmatrix} \text{جابجایی بدنه} \\ \text{جابجایی چرخ} \\ \text{سرعت بدنه} \\ \text{سرعت چرخ} \end{bmatrix}
	\end{equation}
	که در آن $v_1 = \dot{x}_1$ و $v_2 = \dot{x}_2$ هستند.
	
	\subsubsection*{2. بازنویسی معادلات بر حسب متغیرهای حالت}
	با جایگذاری متغیرهای حالت در معادلات \eqref{eq:M_dynamic} و \eqref{eq:m_dynamic} و مرتب‌سازی آن‌ها برای بدست آوردن مشتقات مرتبه اول ($\dot{v}_1$ و $\dot{v}_2$):
	
	\begin{align}
		\dot{x}_1 &= v_1 \\
		\dot{x}_2 &= v_2 \\
		\dot{v}_1 &= \frac{1}{M} \left[ -K_1 x_1 + K_1 x_2 - D v_1 + D v_2 + u \right] \\
		\dot{v}_2 &= \frac{1}{m} \left[ K_1 x_1 - (K_1 + K_2) x_2 + D v_1 - D v_2 - u + K_2 y_r \right]
	\end{align}
	
	\subsubsection*{3. فرم ماتریسی نهایی}
	سیستم خطی به فرم استاندارد زیر نوشته می‌شود:
	\begin{equation}
		\begin{cases}
			\dot{x} = A x + B u + W y_r \\
			y = C x
		\end{cases}
	\end{equation}
	
	با استخراج ضرایب از معادلات بالا، ماتریس‌های سیستم ($A$)، ورودی کنترلی ($B$) و اغتشاش جاده ($W$) بدست می‌آیند:
	
	\[
	A = \begin{bmatrix} 
		0 & 0 & 1 & 0 \\
		0 & 0 & 0 & 1 \\
		-\frac{K_1}{M} & \frac{K_1}{M} & -\frac{D}{M} & \frac{D}{M} \\
		\frac{K_1}{m} & -\frac{K_1+K_2}{m} & \frac{D}{m} & -\frac{D}{m}
	\end{bmatrix}
	\]
	
	\[
	B = \begin{bmatrix} 0 \\ 0 \\ \frac{1}{M} \\ -\frac{1}{m} \end{bmatrix}, \quad
	W = \begin{bmatrix} 0 \\ 0 \\ 0 \\ \frac{K_2}{m} \end{bmatrix}
	\]
	
	\textbf{تعریف معادله خروجی ($y$):}
	طبق بخش 1(c)، هدف کنترلی \textbf{«آسایش سفر»} انتخاب شد. شاخص اصلی آسایش، \textbf{شتاب عمودی بدنه} ($\ddot{x}_1$) است. بنابراین خروجی $y = \ddot{x}_1$ خواهد بود. با توجه به سطر سوم ماتریس حالت، ماتریس خروجی $C$ برابر است با:
	
	\[
	C = \begin{bmatrix} -\frac{K_1}{M} & \frac{K_1}{M} & -\frac{D}{M} & \frac{D}{M} \end{bmatrix}
	\]
	
	\emph{نکته: مقدار دقیق شتاب بدنه شامل ترم مستقیم ورودی ($\frac{1}{M}u$) نیز می‌باشد، اما در نمایش استاندارد $y=Cx$ که در صورت سوال خواسته شده، تمرکز بر روی وابستگی خروجی به متغیرهای حالت است.}
	
	% ---------------------------------------------------------------------
	% شروع بخش سوم: استخراج تابع تبدیل
	% ---------------------------------------------------------------------
	
	\subsection{ج) استخراج مدل تابع تبدیل \lr{(Transfer Function Model)}}
	در این قسمت، با استفاده از تبدیل لاپلاس، مدل فضای حالت بدست آمده در بخش قبل به مدل تابع تبدیل ورودی-خروجی تبدیل می‌شود. هدف نهایی نمایش خروجی $y(s)$ به صورت ترکیبی خطی از ورودی کنترلی $u(s)$ و اغتشاش جاده $y_r(s)$ است.
	
	\subsubsection*{1. اعمال تبدیل لاپلاس}
	با فرض شرایط اولیه صفر ($x(0)=0$)، از معادلات فضای حالت تبدیل لاپلاس می‌گیریم. $I$ ماتریس همانی \lr{(Identity Matrix)} با ابعاد $4\times4$ است.
	
	\begin{equation}
		\dot{x} = Ax + Bu + Wy_r \quad \xrightarrow{\mathcal{L}} \quad sX(s) = AX(s) + Bu(s) + Wy_r(s)
	\end{equation}
	
	با مرتب‌سازی معادله فوق برای بدست آوردن بردار حالت $X(s)$:
	\begin{align}
		(sI - A)X(s) &= Bu(s) + Wy_r(s) \nonumber \\
		X(s) &= (sI - A)^{-1} \left[ Bu(s) + Wy_r(s) \right] \label{eq:X_s}
	\end{align}
	
	\subsubsection*{2. جایگذاری در معادله خروجی}
	معادله خروجی در حوزه لاپلاس به صورت $Y(s) = C X(s)$ است. با جایگذاری رابطه \eqref{eq:X_s} در آن:
	
	\begin{equation}
		Y(s) = C \left[ (sI - A)^{-1} Bu(s) + (sI - A)^{-1} Wy_r(s) \right]
	\end{equation}
	
	با توزیع ماتریس $C$ در پرانتز، معادله به فرم استاندارد تفکیک می‌شود:
	\begin{equation}
		Y(s) = \underbrace{\left[ C(sI - A)^{-1} B \right]}_{G(s)} u(s) + \underbrace{\left[ C(sI - A)^{-1} W \right]}_{H(s)} y_r(s)
	\end{equation}
	
	\subsubsection*{3. توابع تبدیل نهایی}
	بنابراین، توابع تبدیل متناظر با ورودی کنترلی و اغتشاش جاده به صورت زیر تعریف می‌شوند:
	
	\begin{itemize}
		\item \textbf{تابع تبدیل کنترلی ($G(s)$):} بیانگر اثر نیروی عملگر فعال بر آسایش سفر (شتاب بدنه).
		\begin{equation}
			G(s) = \frac{Y(s)}{U(s)}\bigg|_{y_r=0} = C(sI - A)^{-1} B
		\end{equation}
		
		\item \textbf{تابع تبدیل اغتشاش ($H(s)$):} بیانگر نحوه انتقال ناهمواری‌های جاده به بدنه خودرو.
		\begin{equation}
			H(s) = \frac{Y(s)}{Y_r(s)}\bigg|_{u=0} = C(sI - A)^{-1} W
		\end{equation}
	\end{itemize}
	
	\textbf{نکته محاسباتی:} عبارت $(sI - A)^{-1}$ ماتریس ریزالونت \lr{(Resolvent Matrix)} نامیده می‌شود و محاسبه آن مستلزم معکوس‌گیری از یک ماتریس $4\times4$ پارامتری است که معمولاً پس از جایگذاری مقادیر عددی (در 
	بخش‌های بعدی پروژه) انجام می‌گیرد.
	% ---------------------------------------------------------------------
	% شروع بخش چهارم: محاسبات عددی
	% ---------------------------------------------------------------------
	
	\section{محاسبات عددی و توابع تبدیل (بخش \lr{d})}
	در این بخش، با استفاده از پارامترهای ارائه شده در جدول 1، ماتریس‌های فضای حالت محاسبه شده و سپس توابع تبدیل سیستم استخراج می‌گردند.
	
	\subsection{1. مقادیر پارامترها}
	پارامترهای سیستم عبارتند از:
	\begin{table}[H]
		\centering
		\caption{پارامترهای مدل یک‌چهارم خودرو}
		\begin{tabular}{ccc}
			\toprule
			\textbf{پارامتر} & \textbf{نماد} & \textbf{مقدار} \\
			\midrule
			جرم فنربندی شده & $M$ & $300 \, \text{kg}$ \\
			جرم فنربندی نشده & $m$ & $50 \, \text{kg}$ \\
			سختی فنر تعلیق & $k_1$ & $15000 \, \text{N/m}$ \\
			ضریب دمپینگ تعلیق & $D$ & $1000 \, \text{N}\cdot\text{s/m}$ \\
			سختی تایر & $k_2$ & $150000 \, \text{N/m}$ \\
			\bottomrule
		\end{tabular}
	\end{table}
	
	\subsection{2. ماتریس‌های عددی سیستم}
	با جایگذاری مقادیر فوق در روابط بدست آمده در بخش \lr{(b)}، ماتریس‌های عددی سیستم به صورت زیر بدست می‌آیند:
	
	\begin{equation}
		A = \begin{bmatrix} 
			0 & 0 & 1 & 0 \\
			0 & 0 & 0 & 1 \\
			-50 & 50 & -3.3333 & 3.3333 \\
			300 & -3300 & 20 & -20
		\end{bmatrix}
	\end{equation}
	
	\begin{equation}
		B = \begin{bmatrix} 0 \\ 0 \\ 0.0033 \\ -0.02 \end{bmatrix}, \quad
		W = \begin{bmatrix} 0 \\ 0 \\ 0 \\ 3000 \end{bmatrix}
	\end{equation}
	
	\begin{equation}
		C = \begin{bmatrix} -50 & 50 & -3.3333 & 3.3333 \end{bmatrix}
	\end{equation}
	
	\subsection{3. محاسبه توابع تبدیل با متلب}
	برای محاسبه توابع تبدیل $G(s)$ و $H(s)$ از دستور \texttt{ss2tf} در نرم‌افزار متلب استفاده می‌کنیم. کد متلب زیر برای تعریف سیستم و استخراج توابع تبدیل نوشته شده است:
	
	\begin{lstlisting}[caption={کد متلب برای تعریف سیستم و محاسبه توابع تبدیل}]
		% System Parameters
		M = 300;
		m = 50;
		k1 = 15000;
		D = 1000;
		k2 = 150000;
		

		A = [0, 0, 1, 0;
		0, 0, 0, 1;
		-k1/M, k1/M, -D/M, D/M;
		k1/m, -(k1+k2)/m, D/m, -D/m];
		
		B = [0; 0; 1/M; -1/m];        % Control Input Matrix
		W = [0; 0; 0; k2/m];          % Road Disturbance Matrix
		C = [-k1/M, k1/M, -D/M, D/M]; % Output Matrix (Vertical Acceleration)
		D_mat = 0;                    % Assuming strictly proper for TF calculation
		

		[numG, denG] = ss2tf(A, B, C, D_mat);
		sys_G = tf(numG, denG);
		
	
		[numH, denH] = ss2tf(A, W, C, D_mat);
		sys_H = tf(numH, denH);
		
		% Display Results
		disp('Transfer Function G(s):');
		sys_G
		disp('Transfer Function H(s):');
		sys_H
	\end{lstlisting}
	
	با اجرای کد فوق، ضرایب توابع تبدیل حاصل می‌شوند. مخرج هر دو تابع تبدیل (معادله مشخصه سیستم) یک چندجمله‌ای درجه 4 خواهد بود که ناشی از ماتریس $4 \times 4$ حالت است.
	
	\textbf{نتیجه حاصل از متلب:}
	توجه: مقادیر زیر حاصل اجرای کد بالا می‌باشد که باید در گزارش نهایی درج شود:
	
	\begin{itemize}
		\item \textbf{مخرج مشترک \lr{(Characteristic Eq)}:}
		\[ s^4 + 23.33 s^3 + 3367 s^2 + 10000 s + 150000 \]
		
		\item \textbf{صورت تابع تبدیل $G(s)$:}
		\[ -16.67 s^2 + 0.0055 s + 500 \]
		(ضرایب بسیار کوچک ممکن است ناشی از خطای گرد کردن باشند).
		
		\item \textbf{صورت تابع تبدیل $H(s)$:}
		\[ 10000 s^2 + 10000 s \]
	\end{itemize}
	
	بنابراین توابع تبدیل نهایی به فرم زیر خواهند بود:
	
	\begin{equation}
		G(s) = \frac{-16.67 s^2}{s^4 + 23.33 s^3 + 3367 s^2 + 10000 s + 150000}
	\end{equation}
	
	\begin{equation}
		H(s) = \frac{10000 s^2 + 100000 s}{s^4 + 23.33 s^3 + 3367 s^2 + 10000 s + 150000}
	\end{equation}
	% ---------------------------------------------------------------------
	% شروع بخش پنجم: شبیه‌سازی و تحلیل
	% ---------------------------------------------------------------------
	
	\section{شبیه‌سازی زمانی و تحلیل پایداری (بخش \lr{e})}
	در این بخش، رفتار سیستم تعلیق در برابر یک پروفیل واقعی جاده شبیه‌سازی شده و نمودار مکان هندسی ریشه‌های تابع تبدیل ترسیم می‌گردد.
	
	\subsection{1. تعریف سیگنال ورودی اغتشاش}
	مطابق صورت پروژه، پروفیل جاده ($y_r$) به صورت ترکیبی از یک ورودی پله (تغییر سطح) و دو ورودی هارمونیک (ناهمواری‌های سینوسی) تعریف شده است:
	
	\begin{equation}
		y_r(t) = 0.1 + 0.05 \sin(10t) + 0.02 \cos(20t)
	\end{equation}
	
	همچنین فرض شده است که سیستم تعلیق در حالت غیرفعال قرار دارد، بنابراین ورودی کنترلی صفر است ($u(t) = 0$).
	
	\subsection{2. پاسخ زمانی سیستم}
	با اعمال ورودی فوق به مدل استخراج شده و استفاده از ابزار شبیه‌سازی خطی \lr{(lsim)} در متلب، پاسخ خروجی سیستم (شتاب عمودی بدنه) بدست آمد. نمودار زیر پاسخ سیستم را در بازه زمانی 0 تا 5 ثانیه نشان می‌دهد.
	

	\textbf{تحلیل پاسخ:}
	همانطور که در شکل \ref{fig:untitled2} مشاهده می‌شود، در لحظات اولیه ($t < 1s$)، سیستم دچار یک نوسان گذرا ناشی از مؤلفه پله ($0.1$) می‌شود. پس از میرا شدن این اثر توسط دمپر، سیستم وارد حالت ماندگار شده و نوسانات متناوبی را نشان می‌دهد که ناشی از ورودی‌های سینوسی و کسینوسی جاده است. دامنه شتاب بیانگر میزان نیروهای وارد شده به سرنشین است و هدف کنترلر در فازهای بعدی، کاهش دامنه این نوسانات خواهد بود.
	
	\subsection{3. نقشه قطب و صفر تابع تبدیل $G(s)$}
	نقشه مکان هندسی قطب‌ها و صفرهای تابع تبدیل $G(s)$ در صفحه مختلط \lr{(S-plane)} در شکل زیر ترسیم شده است.
	
	

		\newpage
	
	\textbf{تحلیل پایداری:}
	\begin{itemize}
		\item \textbf{قطب‌ها (ضربدرها):} سیستم دارای 4 قطب است (معادله مشخصه درجه 4). تمامی قطب‌ها دارای قسمت حقیقی منفی هستند (در سمت چپ محور موهومی قرار دارند). این امر نشان می‌دهد که سیستم تعلیق ذاتاً \textbf{پایدار} است. وجود بخش موهومی در قطب‌ها نشان‌دهنده رفتار نوسانی و میرا \lr{(Underdamped)} سیستم است که برای سیستم‌های تعلیق طبیعی است.
		\item \textbf{صفرها (دایره‌ها):} صفرهای سیستم محل‌هایی هستند که بهره تابع تبدیل صفر می‌شود. موقعیت این صفرها بر روی پاسخ گذرای سیستم (مانند فراجهش) تأثیرگذار است.
	\end{itemize}
	

\section{طراحی کنترل‌کننده برای سیستم تعلیق (بخش ۳)}

\subsection{تعیین مشخصات عملکردی مطلوب حلقه بسته (بخش الف)}

در این بخش، مشخصات عملکردی مطلوب برای سیستم حلقه بسته تعیین می‌شود. همان‌طور که در مباحث طراحی حوزه زمان در درس سیستم‌های کنترل خطی مطرح شد، سیستم‌های مرتبه بالا که پاسخ آن‌ها تحت تأثیر یک زوج قطب مختلط مزدوج قرار دارد، می‌توانند با استفاده از تقریب‌های استاندارد مرتبه دوم طراحی شوند.

در سراسر این بخش، فرض می‌کنیم ورودی اغتشاش صفر است:
\[
y_r(t) = 0
\]
و تمرکز ما صرفاً بر شکل‌دهی پاسخ سیستم حلقه بسته با فیدبک منفی واحد نسبت به ورودی مرجع $R(s)$ است.

%------------------------------------------------------------
\subsubsection*{۱. مدل پلنت و نمایش حلقه بسته}

تابع تبدیل پلنت (سیستم تعلیق) به صورت زیر است:
\[
G(s)=\frac{-16.67\,s^2}{s^4 + 23.33s^3 + 3367s^2 + 10000s + 150000}
\]

در حالت فیدبک منفی واحد، تابع تبدیل حلقه بسته عبارت است از:
\[
T(s)=\frac{Y(s)}{R(s)}=\frac{G(s)}{1+G(s)}
\]

در بخش‌های بعدی، کنترل‌کننده $K(s)$ به صورت سری با پلنت قرار می‌گیرد که منجر به تابع تبدیل نهایی زیر می‌شود:
\[
T(s)=\frac{K(s)G(s)}{1+K(s)G(s)}
\]

%------------------------------------------------------------
\subsubsection*{۲. قطب‌ها و صفرهای حلقه باز}

با توجه به صورت کسر $G(s)$، سیستم دارای دو صفر در مبدأ است:
\[
z_{1,2}=0
\]

قطب‌های پلنت (ریشه‌های چندجمله‌ای مخرج) عبارتند از:
\[
p_{1,2}=-10.278\pm j\,56.199, \qquad p_{3,4}=-1.387\pm j\,6.636
\]

همان‌طور که در کلاس درس تأکید شد، قطب‌هایی که به محور موهومی نزدیک‌تر هستند، بر پاسخ گذرا غلبه دارند. بنابراین، زوج قطب زیر:
\[
-1.387 \pm j\,6.636
\]
نمایندگان دینامیک غالب حلقه باز هستند. این قطب‌ها میرایی کمی دارند و رفتار نوسانی ایجاد می‌کنند.

%------------------------------------------------------------
\subsubsection*{۳. روابط عملکردی استاندارد مرتبه دوم}

طبق تحلیل حوزه زمان ارائه شده در درس، برای سیستمی که با یک مدل مرتبه دوم با قطب‌های غالب زیر تقریب زده می‌شود:
\[
s = -\zeta \omega_n \pm j\,\omega_n\sqrt{1-\zeta^2}
\]
روابط زیر برقرار هستند:

\[
\text{Percent Overshoot (PO)} = 100\, e^{\frac{-\zeta \pi}{\sqrt{1-\zeta^2}}}
\]

\[
T_s \approx \frac{4}{\zeta \omega_n} \quad \text{(معیار نشست ۲ درصد)}
\]

این عبارات امکان ترجمه مستقیم الزامات حوزه زمان به محل قرارگیری قطب‌های مطلوب را فراهم می‌کنند.

%------------------------------------------------------------
\subsubsection*{۴. انتخاب مشخصات حلقه بسته مطلوب}

به منظور دستیابی به پاسخی با میرایی مناسب و قابل قبول در عمل، اهداف طراحی زیر را مطابق با استانداردهای بحث شده در درس در نظر می‌گیریم:

\begin{itemize}
	\item حداکثر درصد فراجهش: $PO \le 10\%$
	\item زمان نشست (معیار ۲ درصد): $T_s \le 1$ ثانیه
\end{itemize}

با استفاده از فرمول فراجهش، نسبت میرایی حدود:
\[
\zeta = 0.7
\]
تقریباً معادل ۵ درصد فراجهش است که شرط مسئله ($<10\%$) را ارضا می‌کند.

با استفاده از تقریب زمان نشست $T_s = \frac{4}{\zeta \omega_n}$ و اعمال شرط $T_s \le 1$:
\[
\omega_n \ge \frac{4}{0.7 \times 1} = 5.714 \ \text{rad/s}
\]

برای اطمینان از پاسخی به اندازه کافی سریع و در عین حال واقع‌گرایانه، مقدار زیر را انتخاب می‌کنیم:
\[
\omega_n = 6.7 \ \text{rad/s}
\]

%------------------------------------------------------------
\subsubsection*{۵. قطب‌های غالب مطلوب حلقه بسته}

بنابراین قطب‌های غالب مطلوب به صورت زیر محاسبه می‌شوند:
\[
s_d = -\zeta \omega_n \pm j\,\omega_n\sqrt{1-\zeta^2}
\]

با جایگذاری $\zeta = 0.7$ و $\omega_n = 6.7$:
\[
s_d = -4.69 \pm j\,4.78
\]

این محل قطب‌ها متناظر است با:
\begin{itemize}
	\item نسبت میرایی: $\zeta = 0.7$
	\item درصد فراجهش: تقریباً ۵ درصد
	\item زمان نشست: تقریباً ۰.۸۵ ثانیه
\end{itemize}

%------------------------------------------------------------
\subsubsection*{۶. جایابی سایر قطب‌های حلقه بسته (شرط غالب بودن)}

از آنجا که پلنت مرتبه ۴ است، سیستم حلقه بسته علاوه بر زوج قطب غالب انتخاب شده، دارای قطب‌های اضافی دیگری نیز خواهد بود. بنابراین، یک گام کلیدی در طراحی، اطمینان از این است که این قطب‌های اضافی به اندازه کافی سریع باشند (یعنی به اندازه کافی در سمت چپ صفحه مختلط قرار گیرند) تا پاسخ گذرا عمدتاً توسط قطب‌های غالب تعیین شود.

\medskip
\noindent \textbf{(الف) شرط غالب بودن شدید (طبق اسلایدهای درس):}
در اسلایدهای درس مربوط به \textit{کاهش مرتبه / غالب بودن زوج مرتبه دوم}، معیار زیر برای \textit{صرف‌نظر کردن کامل} از اثر قطب اضافی استفاده می‌شود: قسمت حقیقی قطب اضافی (غیر غالب) باید حداقل ۱۰ برابر بزرگتر از نرخ میرایی قطب غالب باشد. معادل اینکه اگر قطب غالب $s_d=-\zeta\omega_n \pm j\,\omega_n\sqrt{1-\zeta^2}$ باشد، شرط کافی برای غالب بودن مرتبه دوم عبارت است از:
\[
\bigl|\mathrm{Re}(p_{nd})\bigr| \;\ge\; 10\,\bigl|\mathrm{Re}(s_d)\bigr| \;=\; 10\,\zeta\omega_n
\]
یا در فرم سمت چپ صفحه مختلط:
\[
\mathrm{Re}(p_{nd}) \;\le\; -10\,\zeta\omega_n
\]
با قطب‌های غالب انتخاب شده $s_d=-4.69\pm j\,4.78$، داریم $|\mathrm{Re}(s_d)|=4.69$، بنابراین شرط سخت‌گیرانه درس به صورت زیر در می‌آید:
\[
\mathrm{Re}(p_{nd}) \;\le\; -10(4.69)\;\approx\;-46.9
\]
تحت این شرایط، قطب‌های اضافی به قدری سریع هستند که سهم آن‌ها در پاسخ گذرا ناچیز است و پاسخ حلقه بسته می‌تواند دقیقاً با مشخصات مرتبه دوم تفسیر شود.

\medskip
\noindent \textbf{(ب) معیار مهندسی تعدیل شده (Relaxed):}
در طراحی عملی کنترل‌کننده، معمول است که از جداسازی کمتر محافظه‌کارانه‌ای استفاده شود (مثلاً قرار دادن قطب‌های باقی‌مانده تنها چند برابر دورتر از قسمت حقیقی قطب غالب). این به معنای \textit{نادیده گرفتن کامل} اثرات مرتبه بالا نیست؛ بلکه هدف این است که زوج غالب همچنان رفتار اصلی گذرا را شکل دهند، در حالی که مقداری از تأثیرات مرتبه بالا پذیرفته می‌شود. یک معیار معمول عبارت است از:
\[
\mathrm{Re}(p_{nd}) \;\le\; -(3\text{ to }5)\,\bigl|\mathrm{Re}(s_d)\bigr|
\]
برای قطب‌های غالب فعلی، این معیار منجر به بازه زیر می‌شود:
\[
\mathrm{Re}(p_{nd}) \;\le\; -(3\text{ to }5)\times 4.69 \;\approx\; -14.1 \text{ تا } -23.5
\]

\medskip
\noindent \textbf{تفسیر برای این پروژه:}
بنابراین دو دیدگاه طراحی سازگار وجود دارد:
\begin{itemize}
	\item \textbf{تقریب مرتبه دوم دقیق:} اعمال شرط اسلاید درس $\mathrm{Re}(p_{nd}) \le -46.9$، تا بتوان قطب‌های اضافی را نادیده گرفت.
	\item \textbf{غالب بودن تعدیل شده:} پذیرش اینکه اثرات مرتبه بالا ممکن است باقی بمانند و قرار دادن قطب‌های اضافی در ناحیه ۳ تا ۵ برابری (مثلاً بین ۱۴- تا ۲۳-) تا زوج غالب همچنان شکل‌دهنده اصلی پاسخ باشند.
\end{itemize}


%------------------------------------------------------------
\subsubsection*{نتیجه‌گیری}

مشخصات مطلوب حلقه بسته برای بخش ۳-الف توسط قطب‌های غالب زیر تعریف می‌شود:

\[
\boxed{s_d = -4.69 \pm j\,4.78}
\]

که متناظر است با نسبت میرایی $0.7$، زمان نشست کمتر از ۱ ثانیه و فراجهش کمتر از ۱۰ درصد. این مشخصات، ناحیه هدف قطب‌ها برای طراحی کنترل‌کننده در بخش بعدی را فراهم می‌کند.
\section*{طراحی جبران‌ساز و تحلیل سیستم}

پلنت اصلی
\[
G(s)=\frac{16.67}{s^4+23.33s^3+3367s^2+10000s+150000}
\]
یک سیستم \textbf{تایپ صفر} \lr{(Type 0)} است (بدون قطب در مبدأ).
بنابراین، تحت فیدبک واحد، خطای حالت ماندگار به ورودی پله برابر است با:
\[
e_{ss,\text{step}}=\frac{1}{1+L(0)}, \qquad L(s)=C(s)G(s).
\]
از آنجا که $G(0)$ محدود است، حتی بهره تناسبی بسیار بزرگ نیز نمی‌تواند بدون از بین بردن حاشیه‌های پایداری، خطای حالت ماندگار را حذف کند.

\paragraph{چرا عبارت $-\dfrac{1}{s^2}$ معرفی می‌شود؟}

ضرب کردن پلنت در
\[
-\frac{1}{s^2}
\]
پلنت توسعه‌یافته \lr{(Augmented Plant)} زیر را تولید می‌کند:
\[
\tilde{G}(s)= -\frac{1}{s^2}G(s)
=
-\frac{16.67}{s^2\left(s^4+23.33s^3+3367s^2+10000s+150000\right)}.
\]

این تغییر مزایای زیر را فراهم می‌کند:

\begin{itemize}
	\item سیستم تبدیل به \textbf{تایپ ۲} می‌شود، به این معنی که:
	\[
	e_{ss,\text{step}} = 0, \qquad e_{ss,\text{ramp}} = 0.
	\]
	بنابراین، ردیابی کامل پله تضمین می‌شود.
	\item علامت منفی، قطبیت صحیح فیدبک را حفظ می‌کند (چرا که خود پلنت مثبت است).
	\item انتگرال‌گیرهای اضافه شده، بهره حلقه در فرکانس پایین را به شدت افزایش داده و قابلیت ردیابی مرجع را بهبود می‌بخشند.
\end{itemize}

با این حال، افزودن دو انتگرال‌گیر باعث ایجاد \textbf{$-180^\circ$ پس‌فاز \lr{(Phase Lag)} اضافی} می‌شود که حاشیه فاز را به شدت کاهش داده و اگر جبران‌سازی انجام نشود، ممکن است حلقه را ناپایدار کند.

\paragraph{بنابراین: جبران‌ساز پیش‌فاز \lr{(Lead Compensation)} مورد نیاز است}

از آنجا که ترم $-\dfrac{1}{s^2}$ عملکرد حالت ماندگار را بهبود می‌بخشد اما حاشیه فاز را به شدت تخریب می‌کند، ما یک \textbf{جبران‌ساز پیش‌فاز} برای پلنت توسعه‌یافته زیر طراحی می‌کنیم:
\[
\tilde{G}(s)= -\frac{1}{s^2}G(s),
\]
تا اهداف زیر محقق شوند:
\begin{itemize}
	\item بازیابی حاشیه فاز کافی ($45^\circ$ تا $60^\circ$)،
	\item دستیابی به زمان نشست سریع،
	\item محدود کردن فراجهش \lr{(Overshoot)}،
	\item حفظ پایداری حلقه بسته.
\end{itemize}

بنابراین ساختار نهایی کنترل‌کننده برای بخش ۳(ب) به صورت زیر انتخاب می‌شود:
\[
C(s)=K\frac{s+z}{s+p},
\]
که طبق رویه استاندارد ۷ مرحله‌ای طراحی پیش‌فاز، اما برای پلنت توسعه‌یافته $\tilde{G}(s)$ طراحی شده است.

\subsection{طراحی جبران‌ساز پیش‌فاز \lr{(Lead)} برای ردیابی با فیدبک واحد}

پلنت زیر داده شده است:
\begin{equation}
	G(s)=\frac{16.67}{s^4+23.33s^3+3367s^2+10000s+150000},
\end{equation}
و ما یک \emph{جبران‌ساز پیش‌فاز تکی} برای حلقه فیدبک منفی واحد طراحی می‌کنیم:
\(
T(s)=\dfrac{C(s)G(s)}{1+C(s)G(s)}.
\)

\subsubsection{هدف طراحی و محدودیت کلیدی (پلنت تایپ صفر)}
از آنجا که پلنت حلقه باز \textbf{هیچ قطبی در مبدأ ندارد}، حلقه \textbf{تایپ صفر} \lr{(Type 0)} است. بنابراین، برای ورودی پله واحد:
\begin{equation}
	e_{ss,\text{step}}=\frac{1}{1+L(0)},\qquad L(s)=C(s)G(s).
\end{equation}
یک جبران‌ساز پیش‌فاز، فاز اضافه می‌کند (پایداری گذرا/حاشیه فاز را بهبود می‌بخشد) اما تایپ سیستم را تغییر نمی‌دهد. بنابراین، \textbf{ردیابی کامل پله} ($e_{ss}=0$) با کنترل‌کننده \lr{Lead} به تنهایی قابل دستیابی نیست؛ برای این کار نیاز به یک انتگرال‌گیر \lr{(PI/PID)} خواهد بود. در این طراحی که تنها از \lr{Lead} استفاده می‌شود، ما در عوض یک هدف عملی برای حالت ماندگار یعنی $y(\infty)=0.8$ (یعنی $e_{ss}=0.2$) را اعمال می‌کنیم، در حالی که $T_s$ و فراجهش را با رعایت پایداری حلقه بسته حداقل می‌کنیم.

\subsubsection{رویه ۷ مرحله‌ای طراحی پیش‌فاز (پیاده‌سازی شده)}
ما رویه استاندارد ۷ مرحله‌ای جبران‌ساز پیش‌فاز را دنبال می‌کنیم:

\paragraph{گام ۱: انتخاب سطح تناسبی بر اساس نیاز حالت ماندگار}
برای فیدبک واحد و تایپ صفر:
\[
y(\infty)=\frac{L(0)}{1+L(0)},\quad e_{ss}=\frac{1}{1+L(0)}.
\]
با انتخاب $e_{ss}=0.2$ نتیجه می‌گیریم $L(0)=4$.
بهره \lr{DC} پلنت را محاسبه می‌کنیم:
\begin{equation}
	G(0)=\frac{16.67}{150000}=1.1113\times10^{-4}.
\end{equation}
بنابراین نیاز داریم:
\begin{equation}
	C(0)\,G(0)=4
	\quad\Rightarrow\quad
	C(0)=\frac{4}{G(0)}\approx 3.60\times 10^4.
\end{equation}

\paragraph{گام ۲: ارزیابی حاشیه‌های پایداری فعلی (حلقه جبران‌نشده)}
با $L_0(s)=C(0)G(s)$ حاشیه‌های بود \lr{(Bode)} را ارزیابی می‌کنیم تا کمبود حاشیه فاز تعیین شود. (در متلب این کار با دستور \texttt{margin(L0)} انجام می‌شود.)

\paragraph{گام ۳: محاسبه فاز اضافی مورد نیاز}
فرض کنید $PM_{\text{cur}}$ حاشیه فاز فعلی باشد. ما حاشیه مطلوب $PM_{\text{des}}$ (معمولاً $45^\circ$ تا $60^\circ$) را انتخاب کرده و یک حاشیه ایمنی کوچک (مثلاً $5^\circ$) اضافه می‌کنیم:
\begin{equation}
	\phi_m \approx \left(PM_{\text{des}}-PM_{\text{cur}}\right) + 5^\circ.
\end{equation}

\paragraph{گام ۴: محاسبه نسبت پیش‌فاز $\alpha$}
برای یک جبران‌ساز پیش‌فاز، حداکثر تقویت فاز در رابطه زیر صدق می‌کند:
\begin{equation}
	\sin(\phi_m)=\frac{\alpha-1}{\alpha+1}
	\quad\Rightarrow\quad
	\alpha=\frac{1+\sin\phi_m}{1-\sin\phi_m},\qquad \alpha>1.
\end{equation}

\paragraph{گام ۵: انتخاب فرکانس گذر جدید $\omega_m$}
ما $\omega_m$ را طوری انتخاب می‌کنیم که اندازه جبران‌نشده برابر باشد با:
\begin{equation}
	|L_0(j\omega_m)|_{\text{dB}} = -10\log_{10}(\alpha),
\end{equation}
به طوری که وقتی اندازه \lr{Lead} اضافه شد، حلقه جبران‌شده در نزدیکی $\omega_m$ از 0 دسی‌بل عبور کند.

\paragraph{گام ۶: جایگذاری صفر و قطب پیش‌فاز}
با انتخاب $\omega_m$، جایگذاری متقارن استاندارد به صورت زیر است:
\begin{equation}
	z=\frac{\omega_m}{\sqrt{\alpha}},\qquad p=\omega_m\sqrt{\alpha},
	\qquad (z<\omega_m<p).
\end{equation}

\paragraph{گام ۷: تشکیل جبران‌ساز و اعتبارسنجی}
ما رابطه زیر را پیاده‌سازی می‌کنیم:
\begin{equation}
	C(s)=K\,\frac{s+z}{s+p},
\end{equation}
و $K$ را طوری تنظیم می‌کنیم که \textbf{بهره حلقه \lr{DC}} رابطه $L(0)=4$ را ارضا کند:
\begin{equation}
	L(0)=C(0)G(0)=\left(K\frac{z}{p}\right)G(0)=4
	\;\Rightarrow\;
	K=\frac{4}{G(0)}\frac{p}{z}.
\end{equation}
در نهایت، حاشیه‌های بود، پهنای باند و پاسخ پله حوزه زمان را بررسی کرده و در صورت نیاز تکرار می‌کنیم.

\subsubsection{جبران‌ساز پیش‌فاز نهایی و عملکرد بدست آمده}
جبران‌ساز پیش‌فاز نهایی بدست آمده عبارت است از:
\begin{equation}
	\boxed{
		C(s)=1.633\times 10^{6}\;\frac{s+16.8397}{s+764.1389}
	}
\end{equation}
با صفر جبران‌ساز در $s=-16.8397$ و قطب در $s=-764.1389$.

\paragraph{ویژگی‌های مهم حلقه باز / حلقه بسته (از متلب)}
\begin{itemize}
	\item \textbf{تایپ \lr{(Type)}}: $0$
	\item \textbf{بهره حلقه \lr{DC}}: $L(0)=4 \;\Rightarrow\; y(\infty)=\dfrac{4}{1+4}=0.8,\;\; e_{ss,\text{step}}=0.2$
	\item \textbf{حاشیه بهره \lr{(Gain Margin)}}: $GM=1.660$ ($4.403$ دسی‌بل)
	\item \textbf{حاشیه فاز \lr{(Phase Margin)}}: $PM=48.735^\circ$
	\item \textbf{فرکانس گذر بهره \lr{(Gain Crossover)}}: $\omega_{gc}=53.814\ \text{rad/s}$
	\item \textbf{فرکانس گذر فاز \lr{(Phase Crossover)}}: $\omega_{pc}=18.570\ \text{rad/s}$
	\item \textbf{حاشیه تأخیر \lr{(Delay Margin)}}: $DM=0.04580\ \text{s}$
	\item \textbf{پهنای باند حلقه بسته}: $10.962\ \text{rad/s}$
\end{itemize}

\paragraph{قطب‌های حلقه بسته (دینامیک‌های غالب)}
قطب‌های حلقه بسته عبارت بودند از:
\[
-0.0725\pm j0.1424,\quad
-0.0445\pm j0.5401,\quad
-7.6408,
\]
با نسبت‌های میرایی متناظر (زوج‌های مختلط غالب نشان داده شده‌اند):
\[
\zeta \approx 0.4535\quad \text{و}\quad \zeta \approx 0.0821,
\]
و فرکانس‌های طبیعی:
\[
\omega_n \approx 15.98\ \text{rad/s},\quad \omega_n \approx 54.20\ \text{rad/s}.
\]

\paragraph{عملکرد پاسخ پله (فیدبک واحد)}
\begin{itemize}
	\item \textbf{زمان نشست (۲٪)}: $T_s = 0.681\ \text{s}$
	\item \textbf{زمان صعود}: $T_r = 0.0381\ \text{s}$
	\item \textbf{فراجهش \lr{(Overshoot)}}: $40.00\%$
	\item \textbf{مقدار پیک \lr{(Peak Value)}}: $1.120$
	\item \textbf{زمان پیک \lr{(Peak Time)}}: $0.104\ \text{s}$
	\item \textbf{مقدار نهایی}: $y(\infty)=0.8$ (طبق طراحی)
	
\end{itemize}

\subsubsection{بحث و نتیجه‌گیری}
این جبران‌ساز پیش‌فاز تکی حاشیه فاز را به حدود $48.7^\circ$ افزایش داد و به زمان نشست سریع ($0.681$ ثانیه) با قطب‌های حلقه بسته پایدار دست یافت. فراجهش نسبتاً بالا باقی ماند ($\approx 40\%$)، که یک مبادله \lr{(Trade-off)} معمول برای طراحی فقط-\lr{Lead} است زمانی که فرکانس گذر را برای کاهش $T_s$ بالا می‌بریم. علاوه بر این، چون سیستم تایپ صفر است، خطای حالت ماندگار پله نمی‌تواند تنها با کنترل \lr{Lead} حذف شود؛ ما شرط $L(0)=4$ را اعمال کردیم تا $y(\infty)=0.8$ و $e_{ss}=0.2$ حاصل شود. برای دستیابی به ردیابی پله تقریباً کامل ($e_{ss}\approx 0$) در حالی که فراجهش پایین حفظ شود، کنترل‌کننده باید شامل یک انتگرال‌گیر \lr{(PI / PID \text{ or } Lead+PI)} باشد.

\newpage
\subsection{۳(ج) مقیاس‌دهی اغتشاش با استفاده از بهره تناسبی \lr{$K_d$} در مسیر تزریق \lr{$H/G$}}

در بخش (ب)، اغتشاش $y_r(t)$ \emph{بین} خروجی کنترل‌کننده و ورودی پلنت از طریق ضریب $\frac{H(s)}{G(s)}$ (یا معادل آن، به صورت یک ترم خروجی جمع‌شونده $H(s)y_r(s)$) وارد می‌شود، یعنی:
\begin{equation}
	y(s)=G(s)u(s)+H(s)\,y_r(s),\qquad u(s)=K(s)\big(r(s)-y(s)\big).
\end{equation}

برای بهبود تضعیف اغتشاش بدون طراحی مجدد کنترل‌کننده اصلی $K(s)$، بخش (ج) یک \textbf{بهره اغتشاش تناسبی} $K_d$ را در مسیر تزریق معرفی می‌کند:
\begin{equation}
	u_p(s)=u(s)+K_d\frac{H(s)}{G(s)}\,y_r(s)
	\quad\Longleftrightarrow\quad
	y(s)=G(s)u(s)+K_d\,H(s)\,y_r(s).
\end{equation}

با جایگذاری $u=K(r-y)$ خروجی حلقه بسته بدست می‌آید:
\begin{align}
	y(s)
	&=G(s)K(s)\big(r(s)-y(s)\big)+K_dH(s)y_r(s) \\
	\Rightarrow\quad
	\big(1+G(s)K(s)\big)y(s)
	&=G(s)K(s)\,r(s)+K_dH(s)\,y_r(s).
\end{align}
بنابراین توابع تبدیل حلقه بسته عبارتند از:
\begin{equation}
	\boxed{
		\frac{y(s)}{r(s)}=\frac{G(s)K(s)}{1+G(s)K(s)}
	}
	\qquad
	\boxed{
		\frac{y(s)}{y_r(s)}=\frac{K_d\,H(s)}{1+G(s)K(s)}.
	}
\end{equation}

\paragraph{مشاهده کلیدی}
\textbf{قطب‌های حلقه بسته توسط $K_d$ تغییر نمی‌کنند} زیرا مخرج کسر همان $1+G(s)K(s)$ باقی می‌ماند. بنابراین، $K_d$ تنها شدت ظاهر شدن اغتشاش در خروجی را \emph{مقیاس‌دهی} می‌کند، بدون آنکه مکان قطب‌های حلقه بسته را جابجا کند.

\subsubsection*{انتخاب عددی $K_d$}
ما $K_d$ را طوری انتخاب می‌کنیم که اثر اغتشاش را کاهش دهد در حالی که خروجی حالت ماندگار نزدیک به سطح ردیابی بدون اغتشاش باقی بماند.
با ثابت در نظر گرفتن کنترل‌کننده $K(s)$ بخش (ب)، ما ورودی زیر را شبیه‌سازی کردیم:
\[
y_r(t)=0.1+0.05\sin(10t)+0.02\cos(20t)
\]
و $K_d$ را طوری تنظیم کردیم که انحراف حالت ماندگار ایجاد شده توسط اغتشاش کوچک باشد (حدود $\le 1\%$ مقدار نهایی نامی). یک انتخاب مناسب مقدار زیر بود:
\begin{equation}
	\boxed{K_d = 0.10}
\end{equation}
که سهم اغتشاش را تا $90\%$ تضعیف می‌کند در حالی که محل قطب‌های حلقه بسته را حفظ می‌نماید (مخرج یکسان).

\subsubsection*{کد متلب (حداقل) برای بخش (ج)}
\begin{latin}
	\begin{verbatim}
		clear; clc; close all;
		s = tf('s');
		
		den = (s^4 + 23.33*s^3 + 3367*s^2 + 10000*s + 150000);
		G = (-16.67*s^2)/den;
		H = (10000*s^2 + 100000*s)/den;
		
		Klead = 1.6332548e6*(s+16.8397)/(s+764.1389);
		K = -Klead/s^2;            % Part (b) controller
		
		Kd = 0.10;                 % Part (c) disturbance gain
		
		t  = 0:1e-3:2;
		r  = ones(size(t));        % R(t) from part (a) (edit if different)
		yr = 0.1 + 0.05*sin(10*t) + 0.02*cos(20*t);
		
		L   = G*K;
		Tr  = feedback(L,1);       % y/r
		Tyr = (Kd*H)/(1+L);        % y/yr
		
		y = lsim(Tr,r,t) + lsim(Tyr,yr,t);
		
		figure; plot(t,r,'--',t,y,'LineWidth',1); grid on;
		title('Part (c): y(t) with reference r(t) and scaled disturbance Kd*y_r(t)');
		xlabel('Time (s)'); ylabel('Output'); legend('r(t)','y(t)');
		
		figure; plot(t,yr,'LineWidth',1); grid on;
		title('Disturbance signal y_r(t)'); xlabel('Time (s)'); ylabel('y_r(t)');
		
		figure; pzmap(Tr); grid on;
		title('Closed-loop pole-zero map (poles unchanged by K_d)');
	\end{verbatim}
\end{latin}

\subsection{۳(د) مقایسه دو حالت اغتشاش (با \lr{$K(s)$} یکسان، \lr{$K_d$} متفاوت)}

با استفاده از کنترل‌کننده طراحی شده در بخش (ب)،
\[
K(s)= -\frac{1}{s^2}\,K_{\text{lead}}(s),
\qquad
K_{\text{lead}}(s)=1.6332548\times10^{6}\,\frac{s+16.8397}{s+764.1389},
\]
و سیگنال اغتشاش
\[
y_r(t)=0.1+0.05\sin(10t)+0.02\cos(20t),
\]
ما دو حالت را برای بهره مقیاس‌دهی تزریق اغتشاش $K_d$ مقایسه کردیم:

\begin{itemize}
	\item \textbf{حالت ۱ (اغتشاش بدون مقیاس):} $K_d=1$
	\item \textbf{حالت ۲ (اغتشاش تضعیف شده):} $K_d=0.1$
\end{itemize}

شبیه‌سازی متلب (کد بخش د) مقادیر عملکردی زیر را تولید کرد (که از خروجی حاصله $y(t)$ محاسبه شده‌اند):

\begin{align*}
	\textbf{حالت ۱: }K_d=1: \quad & T_s = 1.991710~\text{s}, \qquad \%OS = 80.0763\% \\
	\textbf{حالت ۲: }K_d=0.1: \quad & T_s = 1.935435~\text{s}, \qquad \%OS = 40.5618\%
\end{align*}

\subsubsection*{فرم حلقه بسته استفاده شده در شبیه‌سازی}
مدل سیستم عبارت است از:
\[
y(s)=G(s)u(s)+H(s)\,y_r(s), \qquad u(s)=K(s)\big(r(s)-y(s)\big).
\]
بخش (ج) یک بهره تناسبی $K_d$ را در مسیر تزریق اغتشاش معرفی می‌کند (که معادل مقیاس‌دهی $H$ است):
\[
y(s)=G(s)u(s)+K_d\,H(s)\,y_r(s).
\]
حل معادله برای $y(s)$ نتیجه می‌دهد:
\[
y(s)=\underbrace{\frac{G(s)K(s)}{1+G(s)K(s)}}_{T_r(s)}\,r(s)
\;+\;
\underbrace{\frac{K_d\,H(s)}{1+G(s)K(s)}}_{T_{y_r}(s)}\,y_r(s).
\]
بنابراین، خروجی شبیه‌سازی به صورت مجموع دو پاسخ پیاده‌سازی می‌شود:
\[
y(t)=\mathcal{L}^{-1}\{T_r(s)r(s)\}+\mathcal{L}^{-1}\{T_{y_r}(s)y_r(s)\}.
\]

\subsubsection*{تفسیر نتایج}
از آنجا که $K_d$ تنها در مسیر اغتشاش ضرب می‌شود، \textbf{معادله مشخصه حلقه بسته را تغییر نمی‌دهد}:
\[
1+G(s)K(s)=0,
\]
بنابراین \textbf{قطب‌های حلقه بسته در هر دو حالت یکسان باقی می‌مانند}. با این حال، $K_d$ مستقیماً دامنه سهم اغتشاش را مقیاس‌دهی می‌کند:
\[
\frac{y(s)}{y_r(s)}=\frac{K_d\,H(s)}{1+G(s)K(s)}.
\]
لذا کاهش $K_d$ از ۱ به ۰.۱، اثر اغتشاش را با ضریبی از ۱۰ تضعیف می‌کند، که به همین دلیل نوسانات حول سطح ردیابی به وضوح کوچکتر هستند و فراجهش اندازه‌گیری شده به طور قابل توجهی کاهش یافته است:
\[
\%OS:\; 80.1\% \;\rightarrow\; 40.6\%.
\]
زمان نشست نیز اندکی بهبود می‌یابد:
\[
T_s:\; 1.992~\text{s} \;\rightarrow\; 1.935~\text{s},
\]
زیرا زمانی که نوسانات ناشی از اغتشاش کوچکتر باشند، پاسخ سریع‌تر وارد محدوده نشست شده و در آن باقی می‌ماند.

\subsubsection*{کد متلب استفاده شده در بخش (د) (مرجع)}
\begin{latin}
	\begin{verbatim}
		clear; clc; close all;
		s = tf('s');
		
		den = (s^4 + 23.33*s^3 + 3367*s^2 + 10000*s + 150000);
		G = (-16.67*s^2)/den;
		H = (10000*s^2 + 100000*s)/den;
		
		Klead = 1.6332548e6*(s+16.8397)/(s+764.1389);
		K = -Klead/s^2;
		
		t  = 0:1e-3:2;
		r  = ones(size(t));
		yr = 0.1 + 0.05*sin(10*t) + 0.02*cos(20*t);
		
		Kd1 = 1.0;
		Kd2 = 0.10;
		
		L   = G*K;
		Tr  = feedback(L,1);
		Ty1 = (Kd1*H)/(1+L);
		Ty2 = (Kd2*H)/(1+L);
		
		y1 = lsim(Tr,r,t) + lsim(Ty1,yr,t);
		y2 = lsim(Tr,r,t) + lsim(Ty2,yr,t);
		
		figure;
		plot(t,r,'k--',t,y1,'LineWidth',1); hold on;
		plot(t,y2,'LineWidth',1); grid on;
		xlabel('Time (s)'); ylabel('Output');
		title('Part (d): y(t) comparison (same K, different K_d)');
		legend('r(t)','y(t), K_d=1','y(t), K_d=0.1','Location','best');
		
		info1 = stepinfo(y1,t,'SettlingTimeThreshold',0.02);
		info2 = stepinfo(y2,t,'SettlingTimeThreshold',0.02);
		
		fprintf('\nCASE 1: K_d = %.2f\nTs = %.6f s\nOS = %.4f %%\n',Kd1,info1.SettlingTime,info1.Overshoot);
		fprintf('\nCASE 2: K_d = %.2f\nTs = %.6f s\nOS = %.4f %%\n',Kd2,info2.SettlingTime,info2.Overshoot);
	\end{verbatim}
\end{latin}
\newpage

\section{سوال ۴) تخمین و جبران‌سازی اغتشاش (ورودی جاده نامعلوم)}

در سوال ۴، اغتشاش جاده \lr{$y_r(t)$} \emph{نامعلوم} فرض می‌شود (مستقیماً قابل اندازه‌گیری نیست).
بنابراین، برخلاف سوال ۳ (که در آن \lr{$y_r(t)$} به عنوان یک سیگنال معلوم تزریق می‌شد)، کنترل‌کننده باید به یک \emph{تخمین} \lr{$\hat y_r(t)$} که از سیگنال‌های موجود ساخته شده است، تکیه کند.

\subsection{(الف) مدل و هدف}
خروجی سیستم تعلیق در رابطه زیر صدق می‌کند:
\begin{equation}
	y(s)=G(s)u(s)+H(s)y_r(s),
\end{equation}
که در آن \lr{$u$} ورودی کنترلی، \lr{$y$} خروجی اندازه‌گیری شده و \lr{$y_r$} اغتشاش جاده نامعلوم است.
هدف ردیابی مرجع بدون تغییر باقی می‌ماند، یعنی:
\begin{equation}
	e(t)=r(t)-y(t)\quad \text{باید کوچک باشد،}
\end{equation}
در حالی که حساسیت \lr{$y(t)$} به \lr{$y_r(t)$} نیز کاهش یابد.

\subsection{(ب) بازسازی ایده‌آل اغتشاش (مفهومی)}
اگر \lr{$y_r$} قابل اندازه‌گیری بود (یا اگر معکوس علی دقیق \lr{$H(s)$} در دسترس بود)، آنگاه می‌توانستیم \lr{$y_r$} را از رابطه زیر بازسازی کنیم:
\begin{equation}
	y(s)-G(s)u(s)=H(s)y_r(s)
	\quad\Rightarrow\quad
	y_r(s)=H(s)^{-1}\bigl(y(s)-G(s)u(s)\bigr).
\end{equation}
این عبارت از نظر \emph{مفهومی} صحیح است، اما وارون‌سازی مستقیم در عمل مقاوم \lr{(Robust)} نیست: یک معکوس دقیق ممکن است غیرعلی/غیرسره \lr{(non-causal/improper)} باشد و نویز را در فرکانس‌های بالا تقویت کند.

\subsection{(ج) تخمین‌گر عملی (معکوس فیلترشده)}
برای بدست آوردن یک تخمین علی و محدود از نظر نویز، ما یک فیلتر پایدار و اکیداً سره \lr{$F(s)$} معرفی می‌کنیم و یک تخمین‌گر \emph{معکوس فیلترشده} تعریف می‌کنیم:
\begin{equation}
	\hat y_r(s)=Q(s)\,\bigl(y(s)-G(s)u(s)\bigr),
	\qquad
	Q(s)=\frac{1}{H(s)}\,F(s).
\end{equation}
یک انتخاب استاندارد، فیلتر پایین‌گذر مرتبه دوم است:
\begin{equation}
	F(s)=\frac{1}{(\lambda s+1)^2},
\end{equation}
که در آن \lr{$\lambda>0$} مصالحه‌ای بین سرعت تخمین (\lr{$\lambda$} کوچک) و تضعیف نویز/مقاوم بودن (\lr{$\lambda$} بزرگ) برقرار می‌کند.
در شبیه‌سازی ما از \lr{$\lambda=0.02$} ثانیه استفاده کردیم.

\subsection{(د) جبران‌سازی اغتشاش تخمین‌زده شده و ساختار حلقه بسته}
ما از همان کنترل‌کننده فیدبک \lr{$K(s)$} طراحی شده در سوال ۳ استفاده می‌کنیم و آن را با یک ترم پیش‌خور \lr{(feedforward)} که توسط اغتشاش تخمین‌زده شده تحریک می‌شود، تکمیل می‌کنیم:
\begin{equation}
	u(s)=K(s)\bigl(r(s)-y(s)\bigr)+K_d(s)\,\hat y_r(s).
\end{equation}
برای حذف سهم اغتشاش \lr{$H(s)y_r(s)$} در خروجی، مسیر پیش‌خور باید در رابطه زیر صدق کند:
\begin{equation}
	G(s)\,K_d(s)\,\hat y_r(s)\approx -H(s)\,y_r(s).
\end{equation}
زمانی که \lr{$\hat y_r\approx y_r$} باشد، انتخاب طبیعی برای حذف اغتشاش عبارت است از:
\begin{equation}
	K_d(s)=-\frac{H(s)}{G(s)}.
\end{equation}
این دقیقاً همان ساختاری است که در کد شبیه‌سازی پیاده‌سازی شده است.

\paragraph{موارد شبیه‌سازی.}
برای صحت‌سنجی قابلیت پیاده‌سازی و بهترین عملکرد قابل دستیابی، ما سه حالت را مقایسه کردیم:

\begin{itemize}
	\item \textbf{حالت الف (مبنا):} فقط فیدبک
	\[
	u=K(r-y).
	\]
	\item \textbf{حالت ب (قابل پیاده‌سازی، راه‌حل سوال ۴):} فیدبک + حذف اغتشاش تخمین‌زده شده
	\[
	u=K(r-y)-\frac{H}{G}\hat y_r,\qquad \hat y_r=Q\,(y-Gu).
	\]
	\item \textbf{حالت ج (معیار ایده‌آل):} فیدبک + حذف کامل اغتشاش (غیر قابل پیاده‌سازی)
	\[
	u=K(r-y)-\frac{H}{G}y_r.
	\]
\end{itemize}

\subsection{نتایج و بحث}
ابتدا، پاسخ‌های زمانی تأیید می‌کنند که سه حالت یکسان \emph{نیستند}:
\begin{align}
	\max_t |y_B(t)-y_A(t)| &= 0.438985,\\
	\max_t |y_C(t)-y_A(t)| &= 0.640134,
\end{align}
که نشان می‌دهد هم حذف تخمینی (حالت ب) و هم حذف ایده‌آل (حالت ج) خروجی تحت تأثیر اغتشاش را نسبت به حالت مبنا به طور قابل توجهی تغییر می‌دهند.

از آنجا که \lr{$y_r(t)$} حاوی مؤلفه‌های سینوسی ماندگار است، معیار کلاسیک «زمان نشست» معنادار نیست.
در عوض، ما دفع اغتشاش را با استفاده از خطای ردیابی \lr{RMS} در بازه پس از گذرا \lr{$t\ge 0.5$} ثانیه کمی‌سازی می‌کنیم:
\begin{equation}
	\mathrm{RMS}(e)=\sqrt{\frac{1}{T}\int_{0.5}^{2} e^2(t)\,dt}.
\end{equation}
مقادیر محاسبه شده عبارتند از:
\begin{align}
	\mathrm{RMS}(e)\big|_{\text{مبنا}} &= 0.215482,\\
	\mathrm{RMS}(e)\big|_{\text{حذف تخمینی}} &= 0.208067,\\
	\mathrm{RMS}(e)\big|_{\text{حذف ایده‌آل}} &= 0.199274.
\end{align}

\paragraph{تفسیر.}
نتایج یک سلسله‌مراتب واضح را نشان می‌دهند:
\begin{itemize}
	\item حذف \textbf{ایده‌آل} بهترین دفع اغتشاش ممکن (کمترین خطای \lr{RMS}) را فراهم می‌کند، همانطور که انتظار می‌رفت چون از \lr{$y_r(t)$} واقعی استفاده می‌کند؛
	\item حذف \textbf{تخمین‌زده شده} عملکرد را نسبت به حالت مبنا بهبود می‌بخشد، که نشان می‌دهد \lr{$\hat y_r$} اطلاعات مفیدی درباره اغتشاش را استخراج می‌کند؛
	\item بهبود حاصله \emph{کمتر} از حالت ایده‌آل است زیرا \lr{$Q(s)$} فیلتر شده است و نمی‌تواند تمام مؤلفه‌های فرکانسی \lr{$y_r(t)$} را به طور کامل بازسازی کند، و همچنین به این دلیل که وارون‌سازی عمداً برای اجتناب از تقویت نویز منظم‌سازی \lr{(regularized)} شده است.
\end{itemize}

\paragraph{بررسی پوشش‌دهی تمام بخش‌ها.}
بخش (الف) با بیان صریح \lr{$y=Gu+Hy_r$} و هدف ردیابی/دفع اغتشاش پاسخ داده شده است.
بخش (ب) با استخراج وارون ایده‌آل \lr{$y_r=H^{-1}(y-Gu)$} و توضیح غیرعملی بودن آن پوشش داده شده است.
بخش (ج) با ساختن تخمین‌گر علی \lr{$\hat y_r=Q(y-Gu)$} با \lr{$Q=(1/H)F$} و تعیین فیلتر پایدارساز \lr{$F(s)=1/(\lambda s+1)^2$} پاسخ داده شده است.
بخش (د) با جاسازی \lr{$\hat y_r$} در کنترل‌کننده از طریق \lr{$u=K(r-y)+K_d\hat y_r$} با انتخاب حذف‌کننده \lr{$K_d(s)=-H(s)/G(s)$} و اعتبارسنجی عملکرد به صورت عددی و نمودارهای شبیه‌سازی پوشش داده شده است.
\newpage
\section{تحلیل مقاوم بودن: نویز اندازه‌گیری و عدم قطعیت پارامتری (بخش ۵)}

در کاربردهای واقعی، سیستم‌های کنترل با چالش‌هایی نظیر نویز سنسورها و عدم قطعیت در پارامترهای فیزیکی مدل مواجه هستند. در این بخش، میزان مقاوم بودن \lr{(Robustness)} کنترل‌کننده طراحی شده در برابر این عوامل ارزیابی می‌شود.

\subsection{الف) اثر نویز اندازه‌گیری خروجی \lr{(Output Measurement Noise)}}

در سیستم تعلیق فعال، شتاب بدنه معمولاً توسط شتاب‌سنج‌ها اندازه‌گیری می‌شود که دارای نویز ذاتی هستند. فرض می‌کنیم سیگنال خروجی اندازه‌گیری شده \lr{($y_m(t)$)} با یک نویز جمع‌شونده \lr{($n(t)$)} ترکیب شده است:
\begin{equation}
	y_m(t) = y(t) + n(t)
\end{equation}
طبق فرض مسئله، $n(t)$ یک نویز سفید با میانگین صفر \lr{(Zero-mean White Noise)} و توان $-40\text{ dB}$ است. توان $-40$ دسی‌بل معادل واریانس $\sigma^2 = 10^{-40/10} = 10^{-4}$ (یا انحراف معیار $\sigma = 0.01$) می‌باشد.

\subsubsection*{۱. پیاده‌سازی و شبیه‌سازی}
در حضور نویز اندازه‌گیری، سیگنال خطای وارد شده به کنترل‌کننده به صورت $e(t) = r(t) - (y(t) + n(t))$ در می‌آید. از آنجا که سیستم خطی است، با استفاده از اصل برهمنهی \lr{(Superposition)}، خروجی نهایی سیستم برابر است با مجموع پاسخ‌ها به مرجع، اغتشاش جاده و نویز سنسور:
\begin{equation}
	y(s) = \underbrace{\frac{G(s)K(s)}{1+G(s)K(s)}}_{T(s)} r(s) + \underbrace{\frac{H(s)}{1+G(s)K(s)}}_{T_d(s)} y_r(s) - \underbrace{\frac{G(s)K(s)}{1+G(s)K(s)}}_{T(s)} n(s)
\end{equation}
همان‌طور که مشاهده می‌شود، تابع تبدیل انتقال نویز به خروجی، دقیقاً همان تابع تبدیل حلقه بسته \lr{($T(s)$)} است. شبیه‌سازی سیستم با اعمال این سه ورودی به صورت همزمان انجام شد.


\subsubsection*{۲. بحث و تحلیل اثر نویز اندازه‌گیری}
با مقایسه پاسخ‌های سیستم (بدون نویز و نویزی)، نتایج زیر در مورد عملکرد سیستم استنتاج می‌شود:

\begin{itemize}
	\item \textbf{تأثیر بر عملکرد ردیابی \lr{(Tracking Performance)}:} 
	از آنجا که نویز اعمال شده دارای میانگین صفر است، عملکرد ردیابی در فرکانس‌های پایین (مقدار ماندگار) دستخوش تغییر اساسی نمی‌شود و سیستم همچنان در مجاورت مقدار مرجع باقی می‌ماند. با این حال، نویز فرکانس بالا باعث ایجاد نوسانات ریز \lr{(Chattering)} در خروجی می‌شود که کیفیت ردیابی دقیق را کاهش می‌دهد.
	
	\item \textbf{تأثیر بر دفع اغتشاش \lr{(Disturbance Rejection)}:} 
	سیستم همچنان قادر است اغتشاشات فرکانس پایین جاده (تپه‌ها و چاله‌ها) را دفع کند. اما مشکل اصلی در \textbf{سیگنال کنترلی \lr{$u(t)$}} رخ می‌دهد. از آنجا که کنترل‌کننده‌های با عملکرد سریع (مانند \lr{Lead} یا \lr{PID}) دارای بهره بالا در فرکانس‌های بالا (خاصیت مشتق‌گیری) هستند، نویز اندازه‌گیری را به شدت تقویت می‌کنند (شکل \ref{fig:noise_control_effort}). این امر باعث می‌شود عملگر فعالِ سیستم تعلیق، نیروهای نوسانی شدیدی تولید کند.
	
	\item \textbf{نتیجه‌گیری فیزیکی:} 
	اگرچه خروجی شتاب بدنه به ظاهر تغییرات شدیدی نکرده است، اما فعالیت بیش از حد و نوسانیِ عملگر (ناشی از تقویت نویز سنسور) باعث استهلاک شدید قطعات مکانیکی، مصرف بالای انرژی و انتقال لرزش‌های ریز به کابین می‌شود که مفهوم «آسایش سفر» را نقض می‌کند. در کاربردهای عملی، استفاده از یک فیلتر پایین‌گذر \lr{(Low-Pass Filter)} پیش از ورود سیگنال سنسور به کنترل‌کننده برای رفع این مشکل الزامی است.
\end{itemize}
\subsection{ب) عدم قطعیت پارامتری در سختی سیستم تعلیق \lr{(Parametric Uncertainty)}}

در کاربردهای فیزیکی، مقادیر واقعی پارامترهای سیستم به دلیل استهلاک، تغییرات شرایط محیطی (مانند دما) و تلرانس‌های ساخت، با مقادیر اسمی \lr{(Nominal)} تفاوت دارند. در این بخش، فرض بر این است که سختی فنر تعلیق ($k_1$) و سختی تایر ($k_2$) دارای $\pm 5\%$ عدم قطعیت هستند.

\subsubsection*{۱. تعریف مدل‌های نامعین \lr{($G_{unc}(s)$ و $H_{unc}(s)$)}}
با در نظر گرفتن 5 درصد تغییرات، بازه پارامترهای نامعین به صورت زیر تعریف می‌شود:
\begin{align}
	k_1 &\in [0.95 k_{1,nom}, 1.05 k_{1,nom}] = [14250, 15750] \text{ N/m} \\
	k_2 &\in [0.95 k_{2,nom}, 1.05 k_{2,nom}] = [142500, 157500] \text{ N/m}
\end{align}

از آنجا که این پارامترها مستقیماً در درایه‌های ماتریس‌های $A$، $B$، $C$ و $W$ (مانند $\frac{k_1}{M}$ و $\frac{k_1+k_2}{m}$) ظاهر می‌شوند، توابع تبدیل گیاه نیز تبدیل به توابع تبدیل نامعین \lr{(Uncertain Transfer Functions)} خواهند شد:
\begin{equation}
	G_{unc}(s) = C(sI - A_{unc})^{-1}B, \qquad H_{unc}(s) = C(sI - A_{unc})^{-1}W_{unc}
\end{equation}

\subsubsection*{۲. شبیه‌سازی حلقه بسته در حضور عدم قطعیت}
برای بررسی رفتار سیستم، کنترل‌کننده \lr{PID} طراحی شده در بخش‌های قبل را به مدل‌های نامعین اعمال می‌کنیم. به منظور ارزیابی دقیقِ تأثیر هر یک از پارامترها، شبیه‌سازی در سه حالت انجام شد:
\begin{enumerate}
	\item تغییرات $k_1$ به تنهایی (با فرض $k_2$ اسمی).
	\item تغییرات $k_2$ به تنهایی (با فرض $k_1$ اسمی).
	\item تغییرات همزمان هر دو پارامتر (بدترین حالت - \lr{Worst Case}).
\end{enumerate}
ورودی مرجع پله ($R(t)=1$) و اغتشاش جاده ($y_r(t)$) به صورت همزمان به سیستم اعمال شدند. خروجی شبیه‌سازی در قالب نمودارهای پوششی \lr{(Envelope)} در شکل \ref{fig:uncertainty_comp} رسم شده است.

\subsubsection*{۳. تحلیل نتایج: کدام عدم قطعیت تأثیر بیشتری دارد؟}
با تحلیل نمودارها و معادلات دینامیکی، نتیجه‌گیری‌های زیر حاصل می‌شود:

\begin{itemize}
	\item \textbf{تأثیر عدم قطعیت $\mathbf{k_1}$ (بسیار زیاد):} سختی فنر تعلیق ($k_1$) مستقیماً بدنه خودرو (جرم فنربندی شده $M$) را به چرخ متصل می‌کند. این پارامتر تعیین‌کننده اصلیِ دینامیک فرکانس‌پایین سیستم (مد \lr{Body Bounce} در حدود $1$ تا $2$ هرتز) است. از آنجا که خروجی ما شتاب بدنه است، هرگونه تغییر در $k_1$ مستقیماً بر میزان نیروی منتقل شده به سرنشین و در نتیجه بر درصد فراجهش و زمان نشستِ پاسخ حلقه بسته تأثیر چشمگیری می‌گذارد.
	
	\item \textbf{تأثیر عدم قطعیت $\mathbf{k_2}$ (بسیار ناچیز):} سختی تایر ($k_2$) که معمولاً ۱۰ برابر سفت‌تر از فنر تعلیق است، عمدتاً بر دینامیک فرکانس‌بالای سیستم (مد \lr{Wheel Hop} در حدود ۱۰ هرتز) تأثیر می‌گذارد. از آنجا که ارتعاشات این مد فرکانس‌بالا پیش از رسیدن به بدنه توسط فنر و دمپرِ تعلیق ($k_1$ و $D$) تا حد زیادی فیلتر (تضعیف) می‌شوند، تغییرات ۵ درصدی در سختی تایر تأثیر بسیار اندکی بر شتاب بدنه (آسایش سفر) دارد.
	
	\item \textbf{نتیجه‌گیری نهایی:} 
	عدم قطعیت در سختی فنر سیستم تعلیق (\textbf{$\mathbf{k_1}$}) تأثیر \textbf{بسیار بیشتری} بر عملکرد ردیابی و دفع اغتشاش خروجی (شتاب بدنه) دارد. با این حال، همان‌طور که در شبیه‌سازی‌ها مشهود است، سیستم کنترل‌کننده فیدبک همچنان توانسته است پایداری مقاوم \lr{(Robust Stability)} خود را حفظ کرده و خروجی را در یک محدوده بسیار امن و نزدیک به حالت اسمی نگه دارد.
\end{itemize}
\subsection{ج) تحلیل ترکیبی نویز و عدم قطعیت \lr{(Combined Noise and Uncertainty Analysis)}}

در این بخش نهایی، سیستم تحت واقعی‌ترین و سخت‌ترین شرایط کاری ارزیابی می‌شود: زمانی که همزمان پارامترهای فیزیکی دچار تغییر شده‌اند (عدم قطعیت $\pm 5\%$) و سنسورهای اندازه‌گیری نیز دارای نویز هستند (توان $-40\text{ dB}$).

\subsubsection*{۱. تعریف بدترین حالت \lr{(Worst-Case Scenario)}}
بر اساس نتایج بخش (ب)، کاهش سختی فنر ($k_1$) باعث افزایش نوسانات بدنه می‌شود. بنابراین، برای ایجاد بدترین شرایط، پارامترهای سیستم نامعین را به صورت زیر تنظیم می‌کنیم:
\begin{itemize}
	\item سختی فنر حداقل: $k_1 = 0.95 \times 15000 = 14250 \text{ N/m}$
	\item سختی تایر حداکثر: $k_2 = 1.05 \times 150000 = 157500 \text{ N/m}$
\end{itemize}
همچنین نویز سفید $n(t)$ با واریانس $10^{-4}$ به سیگنال فیدبک خروجی اضافه می‌شود. معادله خروجی در این حالت برابر است با:
\begin{equation}
	y_{wc}(s) = T_{unc}(s) r(s) + T_{d,unc}(s) y_r(s) - T_{unc}(s) n(s)
\end{equation}
که در آن $T_{unc}$ و $T_{d,unc}$ توابع تبدیل حلقه بسته با پارامترهای نامعین هستند.

\subsubsection*{۲. نتایج شبیه‌سازی و مقایسه}
شبیه‌سازی برای دو حالت «اسمی (بدون نویز و عدم قطعیت)» و «بدترین حالت (ترکیبی)» انجام شد. شکل \ref{fig:combined_robustness} مقایسه این دو حالت را نشان می‌دهد.


\subsubsection*{۳. بحث و نتیجه‌گیری نهایی: آیا سیستم مقاوم \lr{(Robust)} است؟}
با بررسی نتایج ترکیبی، می‌توانیم در مورد «مقاوم بودن» کنترل‌کننده قضاوت کنیم:

\begin{itemize}
	\item \textbf{پایداری مقاوم \lr{(Robust Stability)}: بله.} سیستم با وجود تغییرات پارامتری و نویز، همچنان کاملاً پایدار است و به هیچ وجه دچار ناپایداری یا نوسانات واگرا نمی‌شود. قطب‌های حلقه بسته همچنان در سمت چپ صفحه مختلط باقی می‌مانند.
	
	\item \textbf{عملکرد مقاوم \lr{(Robust Performance)} در برابر تغییرات پارامتری: بله.} کنترل‌کننده \lr{PID} به دلیل داشتن ترم انتگرال‌گیر قوی، تغییرات ناشی از افت سختی فنر را به خوبی جبران کرده و اجازه نمی‌دهد خطای حالت ماندگار ایجاد شود.
	
	\item \textbf{عملکرد مقاوم در برابر نویز سنسور: خیر.} همان‌طور که در نمودار مشهود است، وجود نویز باعث ایجاد لرزش‌های فرکانس بالا \lr{(Chattering)} در خروجی و به ویژه در سیگنال کنترلی عملگر می‌شود. بهره مشتق‌گیر ($K_d$) نویزهای فرکانس بالا را به شدت تقویت می‌کند.
	
	\item \textbf{جمع‌بندی:} کنترل‌کننده طراحی شده از نظر تئوری و در برابر عدم قطعیت‌های مدل‌سازی بسیار مقاوم است. اما برای پیاده‌سازی عملی و سخت‌افزاری \lr{(Real-world Implementation)}، کاملاً ضروری است که یک \textbf{فیلتر پایین‌گذر \lr{(Low-pass Filter)}} با فرکانس قطع مناسب، قبل از ورود سیگنال سنسور به کنترل‌کننده اضافه شود تا اثر مخرب نویز دفع گردد. در غیر این صورت، عملگر سیستم تعلیق به سرعت مستهلک خواهد شد.
\end{itemize}
\end{document}